\documentclass[10pt,letterpaper]{article}
\usepackage[T1]{fontenc}
\usepackage{spconf,amsmath,url}
\urlstyle{same}
\title{Face--Voice Verification with Protocol-Specific Matching and Kernel Canonical Correlation}
\name{FLAG 2027 local development system}
\address{English training and five-fold identity-held-out evaluation}
\begin{document}
\maketitle

\begin{abstract}
We describe the English-trained system retained after local FLAG development.
It scores a face and a voice with two neural matching ensembles and a
regularized kernel canonical correlation model. The choice of negative
examples during training differs between the standard and same-gender
protocols. The kernel branch learns which training identities have similar
representations in both modalities. Across 100,000 English development
trials, the system achieved 16.604\% equal error rate (EER) on the standard
protocol and 27.216\% on same-gender trials, averaging 21.910\%.
The preceding cross-video ensemble averaged 22.524\% on the same trials.
We report the local results and the archive fingerprint. An official
English--Bangla four-condition EER is unavailable for this system.
\end{abstract}
\begin{keywords}
face--voice verification, kernel canonical correlation analysis,
speaker verification, equal error rate
\end{keywords}

\section{Inputs and training identities}
The FLAG task asks whether a face and a voice belong to the same person.
In the standard protocol, a negative pair can contain any two different
people. In the same-gender protocol, the two people in a negative pair
share a gender. The scorer receives the protocol name and one face--voice
pair~\cite{flag}.

The face input joins frozen CLIP ViT-B/16 and ViT-L/14 image embeddings
into 1,280 components~\cite{clip}. The 2,944-component voice input joins
WavLM, ECAPA-TDNN, ERes2NetV2, WeSpeaker ResNet293-LM, Whisper-small
and a pooled wav2vec 2.0 representation from audEERING's age/gender model.
We use the embedding of the latter model. At inference we do not require
a person's age or gender.

The training pools use English recordings from MAV-Celeb and a screened
subset of VoxCeleb2. For VoxCeleb2 we retained clips of 3--30 seconds
when Whisper-small and VoxLingua107 ECAPA both assigned English a
probability of at least 0.95. This added 1,804 recordings from 88 people.
The language assignments for these clips come from the two frozen
recognizers. The kernel branch used 344--346 training identities per
development fold after removing known aliases. Each of the five outer
folds held out 14 official training identities for evaluation.

\section{Neural matching}
Each network learns a separate 128-dimensional projection for face
and voice. During scoring each projection applies a linear map, frozen
batch normalization, ReLU and unit-length normalization. A Gaussian
model scores the concatenated embeddings under the same-identity and
different-identity hypotheses. The deployment files store the mean and
quadratic-form parameters needed to compute their log-likelihood ratio.

For each protocol, we average two sets of neural scores. One set comes
from the previous English model. Its standard-protocol parameters
exclude the contribution of ERes2NetV2. The other set comes from a
network trained for that protocol. Each set contains three seeds across
five training folds. The local development evaluation uses only the
models belonging to the corresponding held-out fold.

The two protocol-specific networks use the same identity-matching loss.
Within each recording source, the same-gender model pairs each face
identity with other people of the same gender to form negatives. The
standard model draws the same number of negative identities at random
from that source for each face identity. For each protocol we average
the binary cross-entropies on matched pairs and the selected negatives.
This loss and the identity-matching loss each carry weight 0.5. People
with unknown gender still contribute to the identity-matching loss.
Each run uses 750 updates on the English source pool and 200 adaptation
updates on the official training identities.

\section{Cross-modal kernel model}
We unit-normalize each frontend's feature block and give the blocks
equal weight within a modality. For each training identity we remove
repeated file paths and compute separate face and voice centers, giving
each recording equal total weight. Thus a video that supplies many
frames cannot dominate an identity's face center.

For each modality we build a radial basis function kernel using the
training identity centers. The bandwidth is the median squared distance
between distinct identity centers in that modality. We center the
kernel matrix, regularize its eigenspectrum with \(0.01\) times the
matrix trace, and use singular value decomposition to align the face
and voice representations. We retain up to 32 canonical directions
and weight each by the square root of its canonical correlation~\cite{kcca}.

At inference each face or voice is compared with the frozen training
centers and projected into its respective canonical space. The raw
match score is the dot product of the two unit projections. Within each
outer fold we hold out four groups of training identities in turn to
generate scores for logistic calibration. The standard and same-gender
protocols have separate calibration slopes and intercepts. None of these
quantities needs to be fitted on an Evaluation trial.

\section{Scoring a single pair}
Let \(g_0\) and \(g_p\) be the distances from the fixed neural ensemble
and the neural ensemble specialized for protocol \(p\). Let \(k_{fp}\)
be the calibrated kernel log odds from fold \(f\). For the local
deployment archive the fused distance and exported score are
\begin{equation}
 d_p=\frac{g_0+g_p}{2}-\frac{1}{5}\sum_{f=1}^{5} k_{fp},
 \qquad s_p=\frac12+\frac{\arctan(d_p)}{\pi}.
\end{equation}
Smaller scores indicate stronger evidence of a match. Both languages
use the same frozen parameters for a given protocol. The Evaluation
manifest determines the trial order and the destination score file.

\section{Development results}
\begin{table}[!ht]
\centering
\caption{English development EER (\%). Each of five folds contributes
10,000 trials to each protocol, with 5,000 same-identity pairs.}
\label{tab:local}
\smallskip\small\renewcommand{\arraystretch}{1.15}
\begin{tabular}{@{}lrrr@{}}
\hline
Method & Standard & Same-gender & Mean\\ \hline
Cross-video ensemble &16.908&28.140&22.524\\
Protocol-specific models &16.548&27.348&21.948\\
With radial basis kernel &16.604&27.216&\textbf{21.910}\\ \hline
\end{tabular}
\end{table}
The protocol-specific training accounts for 0.576 percentage points
of the 0.614-point mean reduction in Table~\ref{tab:local}. Adding the
kernel changes the mean by a further 0.038 points. A paired bootstrap
over people gives a 95\% interval of \([-1.362,+0.048]\) points for
the full system minus the cross-video ensemble. The kernel configuration
was fixed before reading these results. A linear-kernel comparison
averaged 21.894\%. An earlier system trained with additional languages
averaged 21.790\% on the English trials. We retained the prespecified
radial basis model for the English-trained candidate.

\begin{table}[!ht]
\centering
\caption{EER (\%) on other recordings. Each row compares the same
trials across systems. The two Bangla rows each include seven people.}
\label{tab:transfer}
\smallskip\small\renewcommand{\arraystretch}{1.15}
\begin{tabular}{@{}lrr@{}}
\hline
Corpus and protocol & Cross-video & With kernel\\ \hline
English--Urdu, four conditions &31.038&29.967\\
English--Hindi, standard mean &25.349&25.196\\
Vox Bangla, same gender &22.222&27.778\\
Bangla authors, same gender &46.825&39.683\\ \hline
\end{tabular}
\end{table}
For the bilingual comparisons in Table~\ref{tab:transfer} we refitted
the systems with the corresponding source collections excluded. The
authors' portraits were collected independently of their recordings. The
two Bangla checks moved in opposite directions. Some English and Vox
Bangla development results informed the protocol choices, so they are
reported as development evidence.

\section{Reproducibility}
In the local development repository, \path{reproduce.py} recalculates the
five-fold EERs from saved trial scores and rebuilds the research ZIP
from frozen features. The archived training and scoring entries are
\path{anchor_train.py} and \path{anchor_package.py}. The matcher loads
four neural parameter files and five radial basis models. Replay on 288 trials sampled across
12 local conditions reproduced the stored distances to within
\(2.85\times10^{-14}\). Two packaging runs produced identical archive
bytes for all 2,326 Evaluation trials. The local archive's SHA-256 is
\begin{center}\small\ttfamily
faf1c779c4a86b3dbe5b1a9c09ac5a5c\\
35f3e348c9ffc0f5526afbf5bfcd2ec3
\end{center}

\begin{thebibliography}{9}
\bibitem{flag} FLAG 2027 Organizing Team, \emph{FLAG 2027 Challenge},
\url{https://mavceleb.github.io/dataset/competition.html}.
\bibitem{clip} A. Radford et al., \emph{Learning Transferable Visual Models
From Natural Language Supervision}, Proc. ICML, 2021, pp.~8748--8763.
\bibitem{kcca} D. R. Hardoon, S. Szedmak and J. Shawe-Taylor,
\emph{Canonical Correlation Analysis: An Overview with Application to
Learning Methods}, Neural Computation, vol.~16, 2004, pp.~2639--2664.
\end{thebibliography}
\end{document}
